% Source: 02 - Tue Sep 29 2026.pdf, page 18. Content remains in source order.
\sourcepage{02}{18}
We only require that one certificate exists but do not ask the algorithm to compute it itself.

The certificate is provided to the algorithm externally. ``Computing the certificate'' would imply solving the problem.

A language $L$ is verified in polynomial time if and only if there exists a two input algorithm $V=(x,y)$ that verifies that y is a certificate for x in polynomial time in the size of the input $x$ and the size of the certificate $y$ for all $x\in L$, where y is a certificate for x that must be polynomial in the size of the input $x$.


L is verified in polynomial time iff $\exists V$:
\begin{enumerate}
\item $L=L_V$.
\item $\exists k_1\geq0:\ \forall x\in L\ \exists y:\ V(x,y)=1$ and $|y|=O(|x|^{k_1})$.
\item $\exists k_2\geq0:\ T_V(|x|,|y|)=O((|x|+|y|)^{k_2})$.
\end{enumerate}


The certificate $y$ must also be of polynomial size in the size of the input $x$ otherwise the complexity would be dominated by the size of the certificate and not by the size of the input.

We need SHORT CERTIFICATES!
